\section{Bifibrations and the audit of VIc-C008}

We use ``bifibration'' exclusively for an ordinary Grothendieck fibration which is also an ordinary Grothendieck opfibration. A Thomason bifibration is a different notion. All weak equivalences in this section are in $\W$. In the replacement argument, ``weak fibration'' means universal preservation of $\W$ under strict pullback, the condition of XI Theorem~1(a); for $\W$ this agrees with the weak-fibration tests identified in Part D.

\begin{statement}{E.1}{Proved}{The universal properties required by a strict bifibration}
Let $p:X\to S$ be a Grothendieck fibration, with a cleavage and transport $\alpha^*:X_t\to X_s$ for $\alpha:s\to t$. Then $p$ is an opfibration if and only if each $\alpha^*$ has a left adjoint $\alpha_!:X_s\to X_t$. In that event $B\alpha_!$ and $B\alpha^*$ are mutually inverse homotopy equivalences. Every strict bifibration is both $\W$-proper and $\W$-smooth; the converse for a specified functor is false.
\end{statement}
\begin{proof}
For $x\in X_s$ and $y\in X_t$, cartesian factorization identifies arrows $x\to y$ above $\alpha$ with $\Hom_{X_s}(x,\alpha^*y)$. A cocartesian arrow out of $x$ above $\alpha$ represents this functor of $y$, hence is exactly the datum of $\alpha_!x$ and the adjunction bijection
\[
 \Hom_{X_t}(\alpha_!x,y)\cong\Hom_{X_s}(x,\alpha^*y).
\]
The representing universal properties construct the action on arrows.
They also give the full cocartesian property, not only factorization
against arrows lying over $\alpha$. Indeed, for $\beta:t\to u$ and
$z\in X_u$, the cleavage comparison
$(\beta\alpha)^*\cong\alpha^*\beta^*$ and the adjunction identify
\[
 \Hom^{\beta\alpha}_{X}(x,z)
 \cong \Hom_{X_s}(x,\alpha^*\beta^*z)
 \cong \Hom_{X_t}(\alpha_!x,\beta^*z)
 \cong \Hom^\beta_X(\alpha_!x,z).
\]
Here the superscript specifies the image in $S$, and the composite
bijection is composition with the arrow defined by the adjunction
unit. This proves its universal property for every further base arrow.
Uniqueness then supplies the comparison isomorphisms for composites.
Conversely cocartesian factorization gives the first displayed
adjunction bijections. The unit and counit induce homotopies from the
identity to $B\alpha^*B\alpha_!$ and from $B\alpha_!B\alpha^*$ to the
identity, respectively. Thus no additional weak-equivalence condition
on transports is needed in the question once actual bifibrancy is
known. Cartesian lifts give terminal objects in the incoming lifting
categories; cocartesian lifts give initial objects in the outgoing
lifting categories. The lifting-category tests of Part D therefore
imply smoothness and properness. The finite diamond of D.6 disproves
the converse. In particular the strict condition is stronger, rather
than incomparable, as a property of a given functor.
\end{proof}

\begin{statement}{E.2}{Known}{What Cisinski's factorization actually supplies}
Every functor $p:X\to S$ has a functorial factorization $X\xrightarrow{i}R_SX\xrightarrow{r}S$, with $i\in\W$ and $r$ strongly proper and smooth. If $p$ is a locally constant Grothendieck fibration, $i$ is a weak equivalence on every strict fibre, and $r$ represents the same locally constant object. This statement asserts no ordinary cartesian or cocartesian universal lifting property for $r$.
\end{statement}
\begin{proof}
We opened \citet[Theorem~5.3.14, pp.~230--231]{E_Cisinski2006} before using it. Its cofibrantly generated model structure has the indicated fibrations, so a chosen small-object factorization is functorial. Proposition~5.3.13 gives properness and smoothness, and Corollary~6.4.8 gives local constancy. For the additional fibre assertion, the strict fibres of a Grothendieck fibration with weak-equivalence transports model the homotopy fibres of $Bp$ by Quillen's Theorem B, with arrows reversed. The same holds for the proper and smooth map $r$ by \citet[Corollaries~6.4.8 and 6.4.16]{E_Cisinski2006}. The map $Bi$ is a weak equivalence over $BS$. Naturality of the homotopy-fibre comparison therefore gives weak equivalences $B(X_s)\to B((R_SX)_s)$ for every $s$. This includes empty fibres and components of $S$ outside the image of $p$; a homotopy fibre there is empty on both sides. The specified map, not an unspecified isomorphism of total homotopy types, proves the fibrewise conclusion.
\end{proof}

\begin{statement}{E.3}{Proved}{The exact output of the alternating replacement}
Let $\Phi$ be the comma-category Grothendieck-fibration replacement over $S$, and let $\Psi$ be its opposite. The sequence of VI ms.~192--195,
\[
 X=X_0\longrightarrow X_1=\Phi X\longrightarrow X_2=\Psi\Phi X
 \longrightarrow X_3=\Phi\Psi\Phi X\longrightarrow\cdots,
\]
has a functorial colimit $R^{\mathrm{alt}}_SX$ which is proper and smooth over $S$. Its unit is a total weak equivalence. For locally constant weak-fibration representatives, including locally constant Grothendieck fibrations, the unit is a fibrewise weak equivalence. These assertions do not imply that the colimit is a strict bifibration.
\end{statement}
\begin{proof}
Explicitly, $\Phi_SX$ has objects $(s,x,\alpha:s\to p(x))$ and arrows $(u,v):(s,x,\alpha)\to(s',x',\alpha')$ with $p(v)\alpha=\alpha'u$. Projection to $s$ is a split fibration. The map $i_X$ sends $x$ to $(p(x),x,\id)$; forgetting $s$ retracts it, and the arrows $(\alpha,\id_x)$ give the natural transformation from the identity of $\Phi_SX$ to the resulting composite. Thus $i_X$ is a total weak equivalence. The definition and argument for $\Psi$ reverse all arrows.

Suppose $p$ is a weak fibration. The point inclusion $s\to s\backslash S$ is a weak equivalence over $S$, since $\id_s$ is initial. Its pullback gives $X_s\to s\backslash X=(\Phi_SX)_s$ in $\W$. For $\alpha:s\to t$, precomposition gives a weak equivalence $t\backslash S\to s\backslash S$ over $S$, because both source and target have initial objects. Pulling it back proves that the transports $t\backslash X\to s\backslash X$ are weak equivalences. A Grothendieck fibration with these transports is again a weak fibration: arbitrary strict base change is a Grothendieck fibration with the same transports, and Quillen's Theorem B identifies all strict fibres with homotopy fibres, so pulling back a weak equivalence of bases induces a weak equivalence of total spaces. The argument for $\Psi$ uses $S/s$, with terminal object $\id_s$, and the opposite formulation of Theorem B. These are exactly the stagewise and fibrewise assertions required in the iteration. A locally constant Grothendieck fibration is an allowable initial input by the same application of Theorem B.

For completeness, the colimit step is independent of any universal lifting property. Suppose a filtered system of proper functors is given. An object, a base arrow, an arrow between two lifts, or an equation between such arrows involves finitely many data. The corresponding lifting category in the colimit is consequently the filtered colimit of the lifting categories at sufficiently late stages. The nerve preserves these filtered colimits degreewise. Filtered colimits of weakly contractible simplicial sets are weakly contractible: a map from a finite sphere, a finite homotopy between such maps, and the required finite horn fillers occur at a finite stage after Kan replacement, which commutes with filtered colimits. Thus the outgoing lifting test is retained. The incoming test is dual. The odd stages are smooth, the positive even stages are proper, and both subsequences are cofinal. Their common colimit is therefore proper and smooth.

The same finite-data argument identifies each strict fibre of the colimit with the colimit of strict fibres. Filtered colimits preserve weak equivalences of simplicial sets, so the stagewise comparisons prove the two stated claims about the unit. Naturality and two-out-of-three also show that the replacement preserves fibrewise weak equivalences on the stated subcategory. This justifies the fine-localization comparison with perfect objects, but not a comparison with strict bifibrations. The next statement exhibits the precise universal-property issue.
\end{proof}

\begin{statement}{E.4}{Refuted}{Filtered colimits of strict bifibrations under arbitrary maps}
The claim that a filtered colimit in $\Cat/S$ of strict bifibrations is a strict bifibration is false, already for $S=[1]$.
\end{statement}
\begin{proof}
For $n\geq0$, let $X_n$ be the poset $0<1<\cdots<n<t$, and let $p_n:X_n\to[1]$ send $0,\ldots,n$ to $0$ and $t$ to $1$. Identities are the only endomorphisms, an arrow exists precisely for an inequality, and composition is transitivity. Its fibres are $[n]$ and $e$. The covariant transport $[n]\to e$ has right adjoint selecting $n$, so $p_n$ is a bifibration by E.1. Include $X_n$ in $X_{n+1}$ by the identity on its objects. These are functors over $[1]$. The colimit has fibre $\mathbf N$ over $0$ and $e$ over $1$, with one arrow $k\to t$ for every natural number $k$.

A cartesian lift of $0\to1$ with target $t$ would require a largest natural number. Given a candidate $k\to t$, the arrow $k+1\to t$ cannot factor through it. Hence the colimit is not a fibration. The old cartesian arrow $n\to t$ has already lost its universal property in $X_{n+1}$. Nevertheless every incoming and outgoing lifting category in the colimit is contractible: the only nontrivial incoming one is $\mathbf N$, which has initial object $0$. The colimit remains smooth and proper, exactly as E.3 requires. The code file \texttt{E\_finite\_checks.py} verifies the finite stages and the failure to preserve their cartesian lifts. The unbounded absence of a maximum is the proof of failure in the colimit.
\end{proof}

\begin{statement}{E.5}{Proved}{The classifying space of the contractible-patch category}
Let $Y$ be a small simplicial set and let $\lambda$ be an infinite cardinal in the fixed universe, with $|Y|\leq\lambda$. Let $K_\lambda(Y)$ be the category of nonempty weakly contractible simplicial sets of size at most $\lambda$ mapped to $Y$, with all strict simplicial maps over $Y$. Use fixed subset encodings so this category is small. There is a natural zigzag of weak homotopy equivalences
\[
 BK_\lambda(Y)\ \longleftarrow\ B(\Delta/Y)\ \longrightarrow\ |Y|.
\]
In particular the patch construction in CATHOT D.1 has the correct classifying space; Kan fibrancy is unnecessary for this assertion.
\end{statement}
\begin{proof}
The functor $j_Y:\Delta/Y\to K_\lambda(Y)$ sends a simplex $\sigma:\Delta^n\to Y$ to that same map, and sends a simplex operator to its map of representables. Every $\Delta^n$ is nonempty and contractible and has at most countably many simplices, so it is allowed. For a patch $a:P\to Y$, the comma category $j_Y/(P,a)$ is isomorphic to $\Delta/P$: an object is a simplex $\tau:\Delta^n\to P$, and its composite $a\tau$ specifies the required object of $\Delta/Y$ uniquely. Arrows on both sides are the same commuting triangles of simplices. Its classifying space has the weak homotopy type of $P$, hence is contractible. \citet[\S1, Theorem~A]{E_Quillen1973} implies $Bj_Y\in\W$. The standard category-of-simplices comparison $B(\Delta/Y)\to|Y|$ is natural and a weak equivalence. On a chain of simplex maps ending in $\sigma:\Delta^n\to Y$, it sends barycentric coordinates to the corresponding affine combination of the images of the barycentres of the earlier simplices, then applies $|\sigma|$. Its weak-equivalence property follows by realizing the simplicial replacement of the diagram of standard simplices over $Y$: the map to $B(\Delta/Y)$ collapses its contractible values, and the augmentation to $Y$ contracts the factorization category of each simplex using its identity factorization. This fixes the naturality needed below. For $Y=\varnothing$ both categories are empty. Cardinal coding gives at most $2^\lambda$ objects and arrows and stays inside the fixed universe.
\end{proof}

\begin{statement}{E.6}{Proved}{What patchwise integration represents}
Let $F:S\to\mathbf{sSet}$ be a strict functor taking arrows to weak equivalences, with values bounded in size by $\lambda$. The Grothendieck construction $\int_S K_\lambda(F)$ is an opfibration and represents the covariant local system $F$; equivalently its fibre transport represents the corresponding contravariant local system after inversion in $BS$. If every transition $F(\alpha)$ is a trivial Kan fibration, this opfibration is also a fibration. For a contravariant $F:S^{\op}\to\mathbf{sSet}$ with trivial-Kan-fibration transitions, the fibration obtained by integrating $s\mapsto K_\lambda(F_s)^{\op}$ is a bifibration representing $F$.
\end{statement}
\begin{proof}
Postcomposition makes $K_\lambda(F)$ a strict functor. The fibre over $s$ is $K_\lambda(F_s)$ and the cocartesian transport along $\alpha:s\to t$ is postcomposition with $F(\alpha)$. E.5 is a natural diagramwise comparison, so this fibre diagram is objectwise equivalent to $F$. Integrating the diagram gives its specified local system; the ordinary Grothendieck construction models its homotopy colimit by \citet[Theorem~1.2]{E_Thomason1979}, with the explicit comparison checked in \citet[\S3.1, equation~(9)]{E_Cegarra2020}, after reversing the variance. This is more than a comparison of two unrelated total spaces.

Suppose $f:Y\to Z$ is a trivial Kan fibration. Pullback sends a contractible patch $Q\to Z$ to $Q\times_ZY\to Y$. The projection to $Q$ is a trivial Kan fibration, so the new total simplicial set is nonempty and weakly contractible. The standard slice adjunction restricts to
\[
 f_!:K_\lambda(Y)\rightleftarrows K_\lambda(Z):f^*.
\]
Finite products and pullbacks of sets of size at most the infinite cardinal $\lambda$ retain that bound. If subset codes are used, select isomorphic coded pullbacks; the adjunction is unchanged. E.1, in its covariant formulation, proves bifibrancy. For a contravariant diagram, taking opposite patch categories turns each $f_!$ into a right adjoint, which has the left adjoint required by E.1. The natural homeomorphism $B(C^{\op})\cong BC$, obtained by reversing simplex coordinates, preserves the diagramwise comparison of E.5. This proves the contravariant assertion with its necessary variance. E.10 below supplies the transition hypothesis for a substantial class of bases.
\end{proof}

\begin{statement}{E.7}{Refuted}{Weak-equivalence patch transports need not have adjoints}
Even for a weak equivalence between contractible Kan complexes, postcomposition $K_\lambda(Y)\to K_\lambda(Z)$ can have neither a left nor a right adjoint. Merely enlarging the collection of contractible patches, while retaining strict maps over the target and the same postcomposition functor, does not remedy this example.
\end{statement}
\begin{proof}
Let $G$ be the groupoid on objects $0,1$ with exactly one arrow between every ordered pair. Composition is the unique arrow with the specified endpoints. Both objects are terminal, so $NG$ is a contractible Kan complex. Let $f:\Delta^0\to NG$ select $0$. This is a weak equivalence. Every object in the image of $K_\lambda(f)$ has constant structure map with value $0$. Let $q:\Delta^0\to NG$ select $1$. For every nonempty patch $P$ over $0$, there is no simplicial map $f_!P\to q$ over $NG$, because such a map would equate the distinct zero-simplices $0$ and $1$. There is also no map $q\to f_!P$ over $NG$.

If $f_!$ had a right adjoint $R$, the identity of $R(q)$ would, by the adjunction, produce a map $f_!R(q)\to q$, contradiction. If it had a left adjoint $L$, the identity of $L(q)$ would produce $q\to f_!L(q)$, contradiction. Adding more nonempty contractible domains does not change either zero-simplex argument. This refutes the raw proposed construction, not the existence of a different strict-bifibration representative of its local system. Indeed this particular local system is constant and has an evident product representative. The finite groupoid and the two incompatible point patches are checked in \texttt{E\_finite\_checks.py}.
\end{proof}

\begin{statement}{E.8}{Proved}{All locally constant 1-types have strict bifibration models}
For every small $S$, a locally constant object of $\HOT(S)$ whose values are 1-truncated spaces has a functorial strict-bifibration representative, in the sense of a functorial construction from a strict simplicial presheaf model.
\end{statement}
\begin{proof}
Choose a strict presheaf $F:S^{\op}\to\mathbf{sSet}$ with Kan values representing the object. Apply the fundamental groupoid functor objectwise. The natural maps $F_s\to N\Pi_1F_s$ are weak equivalences because the values are 1-types. Every transition $\Pi_1F_t\to\Pi_1F_s$ is an equivalence of ordinary groupoids, hence has a left adjoint. The ordinary Grothendieck construction $\int_S\Pi_1F$ is consequently a fibration and an opfibration by E.1. Its comparison with $F$ is objectwise and natural, proving the claim in $\HOT(S)$. The construction of the underlying category and projection is functorial in $F$; choosing particular adjoints is unnecessary for the property of being a bifibration. Empty fibres cause no exception: equivalences preserve their emptiness, and the unique functor between two empty categories has itself as both adjoints.
\end{proof}

\begin{statement}{E.9}{Proved}{All bases with 1-truncated classifying space}
If $BS$ is a 1-type, every locally constant object of $\HOT(S)$, with arbitrary space-valued fibres, has a strict-bifibration representative. This includes groupoids, bases with an initial or terminal object, and every weakly contractible base.
\end{statement}
\begin{proof}
Let $G=S[\operatorname{Mor}(S)^{-1}]$ be the ordinary groupoid completion, and let $q:S\to G$ be its canonical functor. The canonical map $BS\to BG$ induces the fundamental-groupoid equivalence; because $BS$ is a 1-type it is a weak equivalence. Restriction along $Bq$ is therefore an equivalence of local-system infinity-categories. Represent the resulting object on $G$ by a strict diagram $H:G^{\op}\to\mathbf{sSet}$ with Kan values, using the ordinary projective model for diagrams of spaces. Since every arrow of $G$ is invertible, every structure map of this strict diagram is a simplicial isomorphism. Choose one cardinal bounding the small family of its values. Then $K_\lambda(H)$ is a diagram of categories with isomorphism transitions. Its Grothendieck construction is a bifibration and models $H$ by E.5. Pulling that bifibration back along the specified $q$ gives a bifibration over $S$, and its diagramwise class is $q^*H$, which is the original local system. This proof uses an equivalence induced by a specified weak-equivalence functor and a diagramwise comparison; it does not infer a fine equivalence from equality of total homotopy types.

The construction is natural after the standard choices of diagram models and replacement functors, but the arbitrary choice of an inverse to restriction is not asserted to be a literal canonical functor on the ordinary homotopy category. Over an initial or terminal base the companion's comparison with a constant diagram gives the more elementary functorial construction directly.
\end{proof}

\begin{statement}{E.10}{Proved}{Bifibration representatives over every poset}
Let $S$ be a small category in which every arrow is a monomorphism, or
every arrow is an epimorphism. Then every locally constant object of
$\HOT(S)$, with arbitrary space-valued fibres, has a small ordinary
strict-bifibration representative. In particular this holds for every
small poset, without a finiteness or dimension restriction. The
construction is functorial on strict simplicial models after a uniform
size bound for their chosen injective fibrant replacements has been
fixed, and has a specified natural diagramwise weak-equivalence
comparison.
\end{statement}
\begin{proof}
First suppose every arrow of $S$ is a monomorphism. Represent the local
system by a strict simplicial presheaf $F$. The injective model
structure on $\mathbf{sSet}^{S^{\op}}$ has objectwise weak equivalences
and objectwise monomorphisms as cofibrations. Its existence and
functorial factorizations follow from
\citet[Proposition A.2.8.2]{FLurieHTT2009}; its compatibility with the
ordinary simplicial enrichment is stated in
\citet[Proposition A.3.3.2 and Remark A.3.3.4]{FLurieHTT2009}.
Choose a functorial injective fibrant replacement $F\to RF$.

For $\alpha:s\to t$, let $h_s=S(-,s)$ and $h_t=S(-,t)$, regarded as
discrete simplicial presheaves. Since $\alpha$ is a monomorphism,
$h_s\to h_t$ is objectwise injective. The simplicial model axiom
therefore makes
\[
 RF_t=\Map(h_t,RF)\longrightarrow\Map(h_s,RF)=RF_s
\tag{E.10a}
\]
a Kan fibration. We spell out its lifting argument. The adjoint of a
horn-lifting problem for (E.10a) has left map
\[
 (h_t\times\Lambda^n_k)\amalg_{h_s\times\Lambda^n_k}
       (h_s\times\Delta^n)
       \longrightarrow h_t\times\Delta^n .
\]
At each object this is a coproduct of identities and horn inclusions,
because $h_s\to h_t$ is injective. It is consequently an objectwise
trivial cofibration, against which the injectively fibrant $RF$ has the
right lifting property. The same argument with the empty presheaf
in place of $h_s$ shows that each $RF_t$ is Kan. Naturality of
$F\to RF$ and two-out-of-three show that (E.10a) is also a weak
equivalence. Thus every transition of $RF$ is a trivial Kan fibration.

Choose an infinite cardinal $\lambda$ bounding all values of $RF$.
The fixed universe is closed under this small family and under the
replacement construction, so such a cardinal exists within it.
E.6 now gives the bifibration
\[
 \int_S\bigl(s\longmapsto K_\lambda(RF_s)^{\op}\bigr)\longrightarrow S.
\tag{E.10b}
\]
The comparison $F\to RF$, the natural simplex-category zigzag of E.5,
and natural reversal of nerves give a specified equivalence of its
fibre diagram with $F$. All maps and compositions are the actual
postcomposition maps of the strict diagram $RF$, so their coherences
are retained. For any small family of input models and maps we can
choose one $\lambda$ bounding all replacement values; with that
choice (E.10b) is functorial. We do not assert that one small cardinal
bounds the values of every small presheaf simultaneously.

If every arrow of $S$ is an epimorphism, use a strict covariant
model of the corresponding local system with transports inverted in
$BS$. The covariant representables $S(t,-)\to S(s,-)$ associated to
$\alpha:s\to t$ are now objectwise injective. The same lifting proof
makes all transitions of an injective fibrant replacement trivial Kan
fibrations, and the covariant construction of E.6 applies.
Every poset satisfies both cancellation hypotheses, proving the
asserted unrestricted poset case.

In particular the six-element poset
$P=\{a_0,a_1\}<\{b_0,b_1\}<\{c_0,c_1\}$ has
$BP\cong S^0*S^0*S^0\cong S^2$, and the path local system classified by
$*\to BP$, with fibre $\Omega S^2$, has the bifibration representative
(E.10b). The program checks its six vertices, twelve edges and eight
triangles. This example is affirmative, not a remaining gap.
\end{proof}

\begin{statement}{E.11}{Open}{The remaining strict-bifibration problem}
For arbitrary small $S$ and arbitrary locally constant $F\in\Fun(S^{\op},\Spaces)$, does there exist an ordinary small strict bifibration $p:X\to S$ whose fibre diagram represents $F$? If so, can one choose it functorially with a specified fibrewise-equivalence comparison from a standard representative? Essential surjectivity outside E.8--E.10, and the stronger general equivalence after fibrewise localization asked in VIc-C008, remain unproved in this audit.
\end{statement}
An explicit remaining test has only one object. Let
\[
 M=\{1\}\amalg\{e_{ij}:i,j\in\{0,1\}\},\qquad
 e_{ij}e_{k\ell}=e_{i\ell},
\tag{E.11a}
\]
with $1$ as identity, and regard $M$ as a one-object category. Every
nonidentity element is idempotent, so none is monic or epic.
The unnumbered lemma of
\citet[pp.~33--34]{E_Fiedorowicz2002} proves $BM\simeq S^2$ for exactly
this monoid. We record its elementary verification. Its universal group
is trivial since all generators are idempotent. Put
$M_i=\{e_{i0},e_{i1}\}$. The right modules $\mathbb Z[M_i]$ are
projective, being the images of idempotent left multiplication on
$\mathbb Z[M]$. There is an exact resolution
\[
 0\longrightarrow\mathbb Z[M_0]\oplus\mathbb Z[M_1]
 \xrightarrow{\mathrm{incl}}\mathbb Z[M]
 \xrightarrow{\,e_{00}-e_{01}\,}\mathbb Z[M_0]
 \xrightarrow{\epsilon}\mathbb Z\longrightarrow0.
\]
The middle map vanishes on all four nonidentity basis elements and
sends $1$ to $e_{00}-e_{01}$, proving exactness directly. Tensoring
with the trivial left module gives
$\mathbb Z^2\xrightarrow{(1,1)}\mathbb Z\xrightarrow{0}\mathbb Z$.
Thus the integral homology is that of $S^2$; simple connectivity,
Hurewicz, and the homological Whitehead theorem give $BM\simeq S^2$.
The code checks the multiplication, module maps, and these integral
chain calculations.

The precise unsettled object is the path local system classified by
$*\to BM$, with fibre $\Omega S^2$. We ask for a strict bifibration over
this one-object, five-arrow category representing that class and having
weakly contractible total category. One object is already minimal for
a nonempty base; we make no claim that five is minimal among possible
numbers of arrows. All poset bases, including the six-element sphere
poset, have been settled by E.10.

The attempted replacements fail at identified steps. E.2 produces homotopical lifting categories, not universal lifts. E.3 does not preserve universal arrows, as E.4 demonstrates. E.7 refutes the assertion that postcomposition on the original patch categories always has the adjoints needed to integrate in both directions; E.10 succeeds by first replacing the whole diagram under its cancellation hypotheses. For (E.11a), the representable maps used in that argument fail to be monomorphisms, so the injective lifting argument does not establish the required transition fibrations. Algebraic Kan fibrations are a possible further modification, but a left adjoint to change of base must also be proved to preserve the contractible-total-space subcategory; preservation does not follow from the existence of algebraic fillers. We do not use that unverified step.

We also searched for and opened a newer free-bifibration construction. \citet[Theorem~1.17 and Propositions~1.23--1.24]{E_ClarkeSchererZeilberger2026} construct a free bifibration on any ordinary functor and a fully faithful unit. Their theorem is a categorical universal property; it does not assert that the unit is a nerve weak equivalence or a fibrewise weak equivalence on locally constant input. That missing homotopical assertion is not supplied here. \citet[\S4]{E_Lamarche2013} explicitly distinguishes his path-based weak equivalences from the standard nerve weak equivalences. We therefore do not convert the free construction into a solution by changing the localizer.

\begin{statement}{E.12}{Proved}{What Chapters XI and XVI settle for this audit}
The lifting conditions of XI \S1 and the factorization arguments of XVI \S4 concern universal preservation of weak equivalences, perfect maps, and locally constant weak fibrations. They do not, by their stated conclusions, settle the ordinary strict-bifibration problem E.11. Two type corrections needed when reading those arguments are $f_s:X'_s\to X_s$ in XI ms.~131 and $p\in\mathcal F$ in the factorization $f=pi$ of XVI Lemma~4.4.1, ms.~41--42.
\end{statement}
\begin{proof}
XI Theorem~1, ms.~127--131, defines a $\Hot$-fibration by stability of weak equivalences under every strict base change, or by the corresponding fibre/comma comparisons after all base changes. No clause specifies a universal cartesian arrow or universal cocartesian arrow. Its Theorem~2 compares $X'\to X$ over $S$, so restriction gives precisely $X'_s\to X_s$. The other typed expression has the wrong domain and codomain for the displayed map.

In XVI Lemma~4.4.1 the output class in a factorization $f=pi$ must contain $p$. Lemma~4.4.3(3) repeats it with this type. The subsequent discussion, ms.~42--43, first factors through a perfect map and then applies $\Phi$ to obtain a locally constant Grothendieck fibration. This second replacement supplies only the fibration direction. The fork on VI ms.~191--192 shows that it can destroy properness, hence cannot silently be declared a bifibration. The comparison of fibration axioms in XVI Theorem~4.4.13 is therefore compatible with E.2 and E.3 and leaves the extra universal-property question distinct.
\end{proof}
